\subsection{Does predicted visibility become reconstructed surface?}
An acquired-mesh audit finds that 7.67\%/7.48\% of G/S expected template reward already lies within 5\,cm of the current native TSDF, averaging prefixes equally. Suppressing only this duplicate reward reduces development-branch mean absolute error, but four complete matched runs remain trajectory- and mesh-identical to their controls. On twelve separately fixed known-geometry objects, the additive mask instead increases gross-area MAE from 0.1563 to 0.1669\,m$^2$. Duplicate removal alone is therefore not a validated improvement in acquisition prediction or exploration.

The forward estimator in Eq.~\ref{eq:acquisition-reward} is then frozen and evaluated on twelve new known-geometry objects, each with revisit, side and around routes. Against the same fine-basis visibility control, MAE decreases from 0.1532 to 0.0143\,m$^2$; one of the 36 dependent route errors increases. Exact hypothetical geometry makes forecast and actual native meshes identical in this probe, so the remaining error reflects the independent surface quadratures. This isolates the acquisition/fusion operator, not inferred-shape accuracy. On all eleven original candidates at the retained action-54 prefix, fitted forecasts improve ranking relative to the matched mean/fine visibility control but slightly worsen absolute error (Table~\ref{tab:acquisition-alignment}). Both select candidate 8 for G and S; the original distribution's S choice is already changed by the mean/basis approximation. No complete exploration run uses this estimator, and no efficiency gain is inferred from the known-geometry result.
\begin{table}[!htbp]
\centering\small
\caption{Prediction versus actual object gross-new area at 5\,cm. Known-geometry rows average three dependent routes on each of twelve new objects; their single error has no G/S prior. Fitted rows use all eleven original candidates at one retained prefix. The last three share conditional-mean shapes and a 0.06\,m quadrature; the original distribution also differs in geometric integration and basis. MAE is in m$^2$; $\rho$ is candidate Spearman correlation. No complete policy uses the acquisition forecast.}
\label{tab:acquisition-alignment}
\begin{tabular}{llrrrr}
\toprule
Geometry & Estimator & MAE (G) & MAE (S) & $\rho_G$ & $\rho_S$ \\
\midrule
Known geometry & Visibility & 0.1532 & -- & -- & -- \\
Known geometry & Acquisition forecast & 0.0143 & -- & -- & -- \\
Fitted geometry & Original distribution & 0.2061 & 0.1961 & 0.927 & 0.891 \\
Fitted geometry & Mean + fine visibility & 0.3107 & 0.3961 & 0.691 & 0.791 \\
Fitted geometry & Mean + native visibility & 0.3070 & 0.3563 & 0.700 & 0.764 \\
Fitted geometry & Mean + acquisition forecast & 0.3172 & 0.3973 & 0.836 & 0.855 \\
\bottomrule
\end{tabular}
\end{table}

